% Standalone section body. Requires amsmath, amssymb, hyperref only.
\section{Hurwitz approximation with a finite witness}
\label{sec:constructive-hurwitz}

The classical constant is retained exactly.  The constructive change is to
replace an unqualified irrational real by a fast rational name together with a
quantitative certificate that it stays away from rationals.  The output is a
rationally checkable strict inequality, not a double-negated existence claim.

\subsection{The input and the exact statement}

Fix primitive recursive functions \(s\mapsto x_s\in\mathbb Q\) and
\(Q\mapsto I(Q)\in\mathbb N\).  In the readable notation below, they name a
real \(x\) and satisfy
\begin{equation}\label{eq:hw-name}
 |x-x_s|\le 2^{-s},\qquad
 |x-p/q|\ge 2^{-I(Q)}
 \quad(p\in\mathbb Z,\ 1\le q\le Q).
\end{equation}
We take \(I\) nondecreasing, replacing it by
\(\max_{1\le j\le Q}I(j)\) when necessary.  This replacement is primitive
recursive.  ``Real'' here is explanatory notation: the first-order formulation
uses rational Cauchy inequalities and, for example, the rational separation
conditions
\[
 |x_s-p/q|+2^{-s}\ge2^{-I(Q)}
 \qquad(1\le q\le Q)
\]
for every indicated index.  Only finitely many instances are used for each
output.  There is no sort of arbitrary real numbers in iPRA.

\paragraph{Constructive Hurwitz theorem.}
For every \(B\in\mathbb N\), a primitive recursive computation from the fixed
functions \(x_s,I\) produces \(p\in\mathbb Z\), \(q,s\in\mathbb N\) such that
\begin{equation}\label{eq:hw-certificate}
 B<q,\qquad \gcd(p,q)=1,\qquad
 5q^4\bigl(|x_s-p/q|+2^{-s}\bigr)^2<1.
\end{equation}
In particular,
\[
 |x-p/q|<\frac1{\sqrt5\,q^2}.
\]
Thus the theorem provides approximants beyond every prescribed denominator,
and therefore an explicitly generated sequence of distinct approximants.
One denominator bound is
\begin{equation}\label{eq:hw-R}
 R_0=1,\qquad R_{j+1}=\max\{R_j,2^{I(R_j)}\},\qquad
 q\le R_{B+3}.
\end{equation}
No useful efficiency is claimed for this bound; its purpose is to make the
primitive recursive character explicit.

As usual for a first-order theory, this is a theorem schema for fixed primitive
recursive definitions of the name and modulus.  It is not an assertion that
iPRA contains a total universal evaluator for all primitive recursive programs.
An arbitrary computable irrational need not supply primitive recursive data
of the required kind.  Dropping the quantitative separation hypothesis would
change the constructive content of the statement.

\subsection{Finite continued fractions and their identities}

Take integers \(a_0\) and \(a_j\ge1\) for \(j\ge1\), and put
\[
 p_{-1}=1,\quad p_0=a_0,\qquad q_{-1}=0,\quad q_0=1,
\]
\[
 p_j=a_jp_{j-1}+p_{j-2},\qquad
 q_j=a_jq_{j-1}+q_{j-2}\quad(j\ge1),
\]
where the recurrence at \(j=1\) uses the displayed initial values in the
usual way (equivalently start with \(p_{-2}=0,q_{-2}=1\)).
Induction on the finite recurrence gives
\begin{equation}\label{eq:hw-det}
 p_jq_{j-1}-p_{j-1}q_j=(-1)^{j-1}\quad(j\ge1).
\end{equation}
Consequently \(p_j,q_j\) are coprime, and so are \(q_j,q_{j-1}\).
Also \(q_1\ge1\), \(q_2\ge2\), and
\begin{equation}\label{eq:hw-growth}
 q_j\ge j\quad(j\ge1),\qquad q_{j+1}\ge q_j+q_{j-1}.
\end{equation}
For the determinant identity, substituting the recurrence changes its sign;
its first value is \(1\).  For the growth bound, \(q_2\ge2\) starts the
induction and the recurrence adds at least one at each later step.

Suppose these are successive digits of the continued fraction of \(x\).
The complete quotient \(t_{j+1}>1\) satisfies
\[
 x=\frac{p_jt_{j+1}+p_{j-1}}{q_jt_{j+1}+q_{j-1}},\qquad
 t_{j+1}>a_{j+1}.
\]
This identity follows by substituting \(t_j=a_j+1/t_{j+1}\) repeatedly;
only a finite number of substitutions is needed.  Subtracting \(p_j/q_j\)
and using \eqref{eq:hw-det} gives
\begin{equation}\label{eq:hw-error}
 \left|x-\frac{p_j}{q_j}\right|
 =\frac1{q_j(q_jt_{j+1}+q_{j-1})}
 <\frac1{q_jq_{j+1}}.
\end{equation}
The sign alternates, so \(x\) lies strictly between consecutive convergents.
For adjacent denominators \(u,v\) and their corresponding numerators \(r,t\),
\begin{equation}\label{eq:hw-bracket}
 \left|x-\frac r u\right|+
 \left|x-\frac t v\right|=\frac1{uv}.
\end{equation}
These are identities about finitely many rational fractional-linear
transformations, with a positive remainder.  They do not require a theorem
about the existence of an infinite continued fraction.

Combining the separation assumption with \eqref{eq:hw-error} yields
\[
 2^{-I(q_j)}\le |x-p_j/q_j|<\frac1{q_jq_{j+1}},
 \qquad q_{j+1}<2^{I(q_j)}.
\]
This proves \eqref{eq:hw-R} by finite induction once the necessary prefix
has been constructed.

\subsection{A rational replacement for the three-convergent argument}

The next lemma is the key to retaining the sharp constant without deciding
comparisons between arbitrary reals.

\paragraph{Denominator lemma.}
Let \(1\le u\le v\), \(w\ge u+v\), and \(\gcd(u,v)=1\).
At least one of the two rational numbers
\[
 S_0=\frac{u^2+v^2}{uv},\qquad
 S_1=\frac{v^2+w^2}{vw}
\]
has square strictly greater than \(5\).  A suitable one is selected by an
integer comparison.

\paragraph{Proof.}
Set \(h=v^2-uv-u^2\).  It cannot vanish.  Indeed, if it did then
\(u\mid v^2\); coprimality implies \(u=1\), and \(v^2=v+1\) has no
positive integral solution: it fails at \(v=1\), while \(v\ge2\) implies
\(v^2\ge v+2\).
If \(h>0\), then
\[
 (u^2+v^2)^2-5u^2v^2
 =(v^2-uv-u^2)(v^2+uv-u^2)>0,
\]
since \(v^2+uv-u^2\ge uv>0\).  Choose \(S_0\).
If \(h<0\), then
\[
 w^2-vw-v^2\ge(u+v)^2-v(u+v)-v^2=-h>0.
\]
The inequality follows either by expansion or from the exact nonnegative
remainder \((w-u-v)(w+u)\).  Applying the preceding factorization to
\(v,w\) gives \(S_1^2>5\).  Choose \(S_1\).
All comparisons in this proof are integer comparisons. \hfill\(\square\)

Apply the lemma to
\[
 (u,v,w)=(q_{B+1},q_{B+2},q_{B+3}).
\]
Every denominator in this triple exceeds \(B\).  Denote the selected pair
again by \(u,v\), and its numerators by \(r,t\).  Put
\begin{equation}\label{eq:hw-margin}
 C=\frac{uv}{u^2+v^2},\qquad d=1-5C^2>0.
\end{equation}
In particular \(0<C\le1/2\), \(0<d<1\), and \(C<1/\sqrt5\).
The familiar classical proof now says that one of the two normalized
errors is at most \(C\).  We instead give a finite selection procedure
that includes a margin certificate.

Let \(Q=\max(u,v)\) and choose a precision \(s\) satisfying
\begin{equation}\label{eq:hw-precision}
 Q^2 2^{-s}\le d/100.
\end{equation}
This need not be found by an unbounded search.  Since
\[
 d=\frac{(u^2+v^2)^2-5u^2v^2}{(u^2+v^2)^2}
 \ge\frac1{(u^2+v^2)^2},
\]
the explicit integer
\[
 s=100Q^2(u^2+v^2)^2
\]
suffices, using \(2^s\ge s\).
Compute the two rational numbers
\[
 A=u^2|x_s-r/u|,\qquad D=v^2|x_s-t/v|
\]
and choose the smaller (with a fixed tie rule).  The triangle inequality
and \eqref{eq:hw-bracket} give
\[
 \frac A{u^2}+\frac D{v^2}
 \le\frac1{uv}+2\cdot2^{-s}.
\]
Therefore their smaller value is at most
\[
 \frac{1/(uv)+2\cdot2^{-s}}{1/u^2+1/v^2}
 =C+2\cdot2^{-s}\frac{u^2v^2}{u^2+v^2}
 \le C+Q^2 2^{-s}.
\]
For the selected numerator and denominator \(p,q\), it follows that
\begin{equation}\label{eq:hw-cert-bound}
 q^2\bigl(|x_s-p/q|+2^{-s}\bigr)
 \le C+2Q^2 2^{-s}\le C+d/50.
\end{equation}
Finally,
\[
 5(C+d/50)^2
 =1-d+Cd/5+d^2/500
 \le1-d+d/10+d/500<1.
\]
Together with \eqref{eq:hw-cert-bound}, this is exactly the rational
certificate \eqref{eq:hw-certificate}.  In particular the proof produces
one of every three indicated convergents and does not use excluded middle
for the real errors.  The constant \(\sqrt5\) disappears entirely from
the certificate checker.

\subsection{Why the continued-fraction prefix can be obtained primitively recursively}

For completeness we give an explicit bounded-precision extraction method.
This is the point at which qualitative irrationality would be insufficient.
Let \(X=\lceil|x_0|\rceil+2\), so \(|x|<X\).
To find \(a_0=\lfloor x\rfloor\), take
\(s_0=I(1)+3\) and compute \(a_0=\lfloor x_{s_0}\rfloor\).
No integer lies within \(2^{-I(1)}\) of \(x\), whereas
\(|x-x_{s_0}|\le2^{-I(1)}/8\).  Thus these two points have the same
integer part.  This also certifies \(0<x-a_0<1\).

Assume the prefix through \(a_{n-1}\) has been obtained, where \(n\ge1\),
and write
\[
 P=p_{n-1},\ Q=q_{n-1},\quad P'=p_{n-2},\ Q'=q_{n-2}.
\]
Here \(Q>0\); for \(n=1\) we use \(P'=1,Q'=0\).
The next complete quotient is
\begin{equation}\label{eq:hw-mobius}
 t_n=\frac{P'-Q'x}{Qx-P}>1.
\end{equation}
Its denominator is nonzero, with the known lower bound
\[
 |Qx-P|\ge D_0:=Q2^{-I(Q)}>0.
\]
Set
\[
 E=\max(1,QX+|P|),\qquad
 U=\left\lceil\frac{|P'|+Q'X}{D_0}\right\rceil+2.
\]
Then \(|Qx-P|\le E\) and \(1<t_n<U-1\).
For every integer \(k\) with \(0\le k\le U\),
\[
 t_n-k=\frac{(P'+kP)-(Q'+kQ)x}{Qx-P}.
\]
Put \(H=Q'+UQ\) and \(L=2^{-I(H)}/E\).
If \(Q'+kQ>0\), irrational separation shows that the absolute value
of the numerator is at least \(2^{-I(H)}\).
The only possible zero-denominator coefficient is \(Q'=k=0\), when
\(|P'|=1\) by the determinant identity, giving the same lower bound.
Consequently
\begin{equation}\label{eq:hw-floor-gap}
 |t_n-k|\ge L>0\qquad(0\le k\le U).
\end{equation}

Choose a rational precision bound
\[
 \delta=\min\{D_0/(4Q),\ LD_0^2/8\}>0,
\]
and choose \(s\) with \(2^{-s}\le\delta\) by a bounded rational computation.
For instance, if \(\delta=a/b>0\) in integer form, \(s=b+1\) suffices;
no reduction of the fraction is necessary.  Put \(y=x_s\).
The first bound gives \(|Qy-P|\ge3D_0/4\).  The determinant identity gives
the exact difference formula
\[
 \left|
 \frac{P'-Q'y}{Qy-P}-\frac{P'-Q'x}{Qx-P}
 \right|
 =\frac{|x-y|}{|Qy-P|\,|Qx-P|}
 \le\frac{4\delta}{3D_0^2}\le L/6.
\]
Here $L\le1$, and $1<t_n<U-1$ ensures that only the integers in
$[0,U]$ can matter.  By \eqref{eq:hw-floor-gap}, the two quotients have the same integer part.
Thus exact rational division and flooring of the quotient at \(y\) gives
\(a_n\), and the associated inequalities certify
\(a_n<t_n<a_n+1\).  Set \(t_{n+1}=1/(t_n-a_n)\) in the explanatory
notation and update the integer convergents by their recurrences.

Iterating this finite procedure through \(a_{B+3}\) supplies the prefix
used above.  Every operation is integer or rational arithmetic, a bounded
maximum, fixed-length primitive recursion, or evaluation of the fixed
primitive recursive input functions.  Each claimed real inequality has a
strict rational margin and can be replaced by inequalities between sufficiently
accurate rational names.  The loop invariants consist of determinant,
positivity, and rational enclosure certificates; they are decidable finite
conditions.  This is the promised path to a quantifier-free induction proof.

An alternative interface accepts a primitive recursive sequence of positive
continued-fraction digits directly.  Its convergents give a primitive recursive
Cauchy name (the denominators grow at least as Fibonacci numbers).  A
quantitative irrationality witness also follows: for \(1\le q\le B\), choose
an explicitly indexed convergent with \(q_n>B\) and \(q_{n+1}\ge2B\).
Then, unless the rational equals that convergent (which is impossible in
lowest terms),
\[
 |x-p/q|\ge\frac1{Bq_n}-\frac1{q_nq_{n+1}}
 \ge\frac1{2Bq_n}.
\]
A later index such as \(n=2B+2\) suffices by \eqref{eq:hw-growth}.
This explains the equivalence of the two useful input interfaces without
postulating an unbounded continued-fraction oracle.
For example, $I(B)=2Bq_{2B+2}$, for $B\ge1$, is an explicit sufficient
modulus for the digit interface, since $2^r\ge r$; it may be replaced
by its bounded running maximum.

\subsection{An entirely finite certificate interface}

The use of $x$ above is eliminable at the main arithmetic step.
Supply the triple of integer convergents and a rational interval
$[\ell,h]$ contained in both adjacent-convergent brackets.  Its midpoint
$y$ and half-width $e$ satisfy $|z-y|\le e$ for every $z$ in that interval.
For the pair selected by the denominator lemma, require the rational
test $Q^2e\le d/100$.  The same weighted-error calculation, now entirely
over rational endpoints, selects $p/q$ and proves
\[
 5q^4\max\{|\ell-p/q|,|h-p/q|\}^2<1.
\]
The convexity inequality for absolute value on an interval follows by
the three rational cases $p/q\le\ell$, $\ell<p/q<h$, and $h\le p/q$.
Consequently this is a finite rational theorem with no real parameters.

In the name interface, refine the final precision further so that
$2^{-s}\le2^{-I(Q_*)-3}$, where $Q_*$ is the largest denominator of
the triple.  Each convergent endpoint stays at least $2^{-I(Q_*)}$
from the named value.  Thus $[x_s-2^{-s},x_s+2^{-s}]$ lies inside both
adjacent-convergent brackets, and the enclosure checks are finite rational
checks.  The added precision is a primitive recursive maximum of the
two prescribed bounds and does not change the denominator bound.

The same enclosure discipline applies during digit extraction.  With
the displayed choice of $\delta$, every rational $z$ in
$[y-2^{-s},y+2^{-s}]$ has $|z-x|\le2\delta$ in the explanatory
notation, hence $|Qz-P|\ge D_0/2$.  The fractional-linear value changes
by at most $4\delta/D_0^2\le L/2$, so the rational image interval
crosses no integer.  Its two endpoint floors certify the next digit.
Coherence of the name supplies compatible later refinements.  This
gives the concrete finite-enclosure invariant to be implemented in
the object-theory formalization.

\subsection{Optimality, with an explicit finite denominator bound}

Let \(\phi=(1+\sqrt5)/2\) and \(\phi'=(1-\sqrt5)/2\).
These numbers have primitive recursive rational names by bisection of
\(z^2-5\) on \([2,3]\), using a prescribed finite number of bisections.
For every \(p\in\mathbb Z,q\ge1\),
\begin{equation}\label{eq:hw-norm}
 N=p^2-pq-q^2\ne0,\qquad
 1\le|N|=q^2|p/q-\phi|\,|p/q-\phi'|.
\end{equation}
To verify that \(N\ne0\), divide \(p,q\) by their greatest common divisor.
If \(N=0\), coprimality would force the reduced denominator to be one,
and an integer \(z\) would satisfy \(z(z-1)=1\).  The product is zero
at \(z=0,1\), and at least two for every other integer, a contradiction.
This proof uses only integer arithmetic.

The same identity supplies the required irrationality modulus.  If
\(0\le p/q\le3\), then \(|p/q-\phi'|<4\), whence
\(|\phi-p/q|>1/(4q^2)\).  Outside that rationally decidable range the
distance is greater than one.  Hence
\[
 |\phi-p/q|\ge\frac1{4Q^2}\ge2^{-4Q^2}
 \qquad(1\le q\le Q),
\]
so \(I(Q)=4Q^2\) is a valid elementary modulus.

Now let \(c>0\) be rational with \(5c^2<1\), and put
\(d=1-5c^2>0\).  Suppose
\[
 |\phi-p/q|<c/q^2.
\]
Since \(|\phi-\phi'|=\sqrt5\), \eqref{eq:hw-norm} gives
\[
 1<c\sqrt5+c^2/q^2.
\]
But
\[
 1-c\sqrt5=\frac{1-5c^2}{1+c\sqrt5}>d/2.
\]
Therefore
\begin{equation}\label{eq:hw-optimal-bound}
 q^2<\frac{2c^2}{1-5c^2}.
\end{equation}
Only finitely many denominators are possible; for each denominator the
approximation inequality also bounds the numerator.  Every proposed rational
coefficient strictly smaller than \(1/\sqrt5\) therefore fails to give
approximants beyond every denominator for this explicitly presented irrational.
The argument proves optimality constructively, with the explicit bound
\eqref{eq:hw-optimal-bound}, rather than merely by a limiting statement about
Fibonacci ratios.

\subsection{Lean implementation plan for this part}

The arithmetic core should be formalized over \(\mathbb Z\) and a rational
number representation, not by importing a classical real-number Hurwitz
theorem.  A useful dependency order is:
\begin{enumerate}
\item finite signed-integer and rational arithmetic, powers of two, bounded
      maxima, and exact rational comparison;
\item convergent recurrences and the determinant, coprimality, and growth
      identities;
\item the denominator lemma, using its two displayed polynomial
      factorizations and the decidable sign of \(h\);
\item a rational bracketing lemma: a point in an interval of length \(1/(uv)\)
      and its \(\varepsilon\)-approximation satisfy the weighted-error bound;
\item the selector and rational certificate \eqref{eq:hw-certificate}, with
      \(C,d,s\) exactly as above;
\item bounded-precision prefix extraction from a name and a separation
      modulus, maintaining finite interval and determinant certificates;
\item the integer norm argument for the golden ratio and its bounded
      bisection name.
\end{enumerate}
The main internal statement can use the rational inequality
\(5q^4(|x_s-p/q|+2^{-s})^2<1\) directly.  A separate interpretation lemma
then reads this as the familiar real inequality.  Lean's ambient logic does
not by itself certify an iPRA derivation: the extraction functions need
primitive-recursion certificates, and their arithmetic proofs need to be
compiled into the chosen iPRA syntax or checked against its induction
restrictions.  The finite identities and invariants above make those
obligations explicit.

\paragraph{Classical reference.}
For the classical three-convergent theorem and the sharp constant, see
Michael Filaseta, \emph{Math 784: Algebraic Number Theory, Instructor's Notes},
Theorem 25, printed pages 20--21,
\url{https://people.math.sc.edu/filaseta/gradcourses/TheMath784Notes.pdf}.
The rational pair selector, explicit margins, certificate interface, and
bounded-precision extraction above give the constructive implementation
needed here; the classical statement alone would not supply it.
